\documentclass[a4paper, final]{article}
%\usepackage{literat} % Нормальные шрифты
\usepackage[14pt]{extsizes} % для того чтобы задать нестандартный 14-ый размер шрифта
\usepackage{tabularx}
\usepackage[T2A]{fontenc}
\usepackage[utf8]{inputenc}
\usepackage[russian]{babel}
\usepackage{amsmath}
\usepackage[left=25mm, top=20mm, right=20mm, bottom=20mm, footskip=10mm]{geometry}
\usepackage{ragged2e} %для растягивания по ширине
\usepackage{setspace} %для межстрочного интервала
\usepackage{moreverb} %для работы с листингами
\usepackage{indentfirst} % для абзацного отступа
\usepackage{moreverb} %для печати в листинге исходного кода программ
\usepackage{pdfpages} %для вставки других pdf файлов
\usepackage{tikz}
\usepackage{graphicx}
\usepackage{afterpage}
\usepackage{longtable}
\usepackage{float}

% \usepackage[paper=A4,DIV=12]{typearea}
\usepackage{pdflscape}
% \usepackage{lscape}

\usepackage{array}
\usepackage{multirow}

\renewcommand\verbatimtabsize{4\relax}
\renewcommand\listingoffset{0.2em} %отступ от номеров строк в листинге
\renewcommand{\arraystretch}{1.4} % изменяю высоту строки в таблице
\usepackage[font=small, singlelinecheck=false, justification=centering, format=plain, labelsep=period]{caption} %для настройки заголовка таблицы
\usepackage{listings} %листинги
\usepackage{xcolor} % цвета
\usepackage{hyperref}% для гиперссылок
\usepackage{enumitem} %для перечислений

\setlist[enumerate,itemize]{leftmargin=1.2cm} %отступ в перечислениях

\hypersetup{colorlinks,
	allcolors=[RGB]{010 090 200}} %красивые гиперссылки (не красные)

% подгружаемые языки — подробнее в документации listings
\lstloadlanguages{Haskell}
\lstset{tabsize=2,
	breaklines,
	basicstyle=\footnotesize,
	columns=fullflexible,
	flexiblecolumns,
	numbers=left,
	numberstyle={\footnotesize},
	keywordstyle=\color{blue},
	inputencoding=utf8,
	extendedchars=true,
	literate={а}{{\cyra}}1 {б}{{\cyrb}}1 {в}{{\cyrv}}1 {г}{{\cyrg}}1 {д}{{\cyrd}}1 {е}{{\cyre}}1 {ё}{{\"e}}1 {ж}{{\cyrzh}}1 {з}{{\cyrz}}1 {и}{{\cyri}}1 {й}{{\cyrishrt}}1 {к}{{\cyrk}}1 {л}{{\cyrl}}1 {м}{{\cyrm}}1 {н}{{\cyrn}}1 {о}{{\cyro}}1 {п}{{\cyrp}}1 {р}{{\cyrr}}1 {с}{{\cyrs}}1 {т}{{\cyrt}}1 {у}{{\cyru}}1 {ф}{{\cyrf}}1 {х}{{\cyrh}}1 {ц}{{\cyrc}}1 {ч}{{\cyrch}}1 {ш}{{\cyrsh}}1 {щ}{{\cyrshch}}1 {ъ}{{\cyrhard}}1 {ы}{{\cyry}}1 {ь}{{\cyrsoft}}1 {э}{{\cyre}}1 {ю}{{\cyryu}}1 {я}{{\cyrya}}1 {А}{{\cyrA}}1 {Б}{{\cyrB}}1 {В}{{\cyrV}}1 {Г}{{\cyrG}}1 {Д}{{\cyrD}}1 {Е}{{\cyrE}}1 {Ё}{{\"E}}1 {Ж}{{\cyrZH}}1 {З}{{\cyrZ}}1 {И}{{\cyrI}}1 {Й}{{\cyrISHRT}}1 {К}{{\cyrK}}1 {Л}{{\cyrL}}1 {М}{{\cyrM}}1 {Н}{{\cyrN}}1 {О}{{\cyrO}}1 {П}{{\cyrP}}1 {Р}{{\cyrR}}1 {С}{{\cyrS}}1 {Т}{{\cyrT}}1 {У}{{\cyrU}}1 {Ф}{{\cyrF}}1 {Х}{{\cyrH}}1 {Ц}{{\cyrC}}1 {Ч}{{\cyrCH}}1 {Ш}{{\cyrSH}}1 {Щ}{{\cyrSHCH}}1 {Ъ}{{\cyrHARD}}1 {Ы}{{\cyrY}}1 {Ь}{{\cyrSOFT}}1 {Э}{{\cyrE}}1 {Ю}{{\cyrYU}}1 {Я}{{\cyrYA}}1
}

\textheight=24cm % высота текста
\textwidth=16cm % ширина текста
\oddsidemargin=0pt % отступ от левого края
\topmargin=-1.5cm % отступ от верхнего края
\parindent=24pt % абзацный отступ
\parskip=5pt % интервал между абзацами
\tolerance=2000 % терпимость к "жидким" строкам
\flushbottom % выравнивание высоты страниц

% Настройка листингов
\lstset{
	language=Haskell,
	extendedchars=\true,
	inputencoding=utf8,
	keepspaces=true,
	captionpos=t,
}

\begin{document} % начало документа
	
	% НАЧАЛО ТИТУЛЬНОГО ЛИСТА
	\begin{center}
		\hfill \break
		\hfill \break
		\normalsize{МИНИСТЕРСТВО НАУКИ И ВЫСШЕГО ОБРАЗОВАНИЯ РОССИЙСКОЙ ФЕДЕРАЦИИ\\
			федеральное государственное автономное образовательное учреждение высшего образования «Санкт-Петербургский политехнический университет Петра Великого»\\[10pt]}
		\normalsize{Институт компьютерных наук и кибербезопасности}\\[10pt] 
		\normalsize{Высшая школа технологий искусственного интеллекта}\\[10pt] 
		\normalsize{Направление: 02.03.01 <<Математика и компьютерные науки>>}\\
		
		\hfill \break
		\hfill \break
		\hfill \break
		\hfill \break
		\large{Отчет по курсовой работе}\\
		\large{по дисциплине}\\
		\large{<<Функциональное программирование>>}\\
		\large{Вариант 16}\\
		\large{Тема: Модификатор изображений}\\
		\hfill \break
		
	\end{center}
	
	\small{ 
		\begin{tabular}{lrrl}
			\!\!\!Студент, & \hspace{2cm} & & \\
			\!\!\!группы 5130201/30102 & \hspace{2cm} & \underline{\hspace{3cm}} & Санько В. В. \\\\
			\!\!\!Преподаватель,\\ \hspace{-5pt}к. т. н., доц. & \hspace{2cm} &  \underline{\hspace{3cm}} &  Моторин Д. Е. \\\\
			&&\hspace{4cm}
		\end{tabular}
		\begin{flushright}
			<<\underline{\hspace{1cm}}>>\underline{\hspace{2.5cm}} 2025г.
		\end{flushright}
	}
	
	\hfill \break
	\begin{center} \small{Санкт-Петербург, 2025} \end{center}
	\thispagestyle{empty} 
	
	% КОНЕЦ ТИТУЛЬНОГО ЛИСТА
	\newpage
	
	\tableofcontents
		\newpage
	\section {Постановка задачи}
	
	Целью курсовой работы является разработка программного комплекса для обработки растровых изображений на языке Haskell. Работа разделена на две части:
	
	\begin{enumerate}
		\item \textbf{Часть 1:} Воспроизведение решения задачи обработки изображений, описанной в главах 12 и 13 книги Филиппа Хагенлохера «Haskell from the Very Beginning».
		Требуется:
		\begin{itemize}
			\item Реализовать парсер формата PNM (PPM/PGM) с использованием комбинаторных библиотек.
			\item Использовать обобщенные алгебраические типы данных (GADT) для типизации изображений.
			\item Реализовать параллельную обработку данных.
		\end{itemize}
		
		\item \textbf{Часть 2:} Разработка собственного синтаксического анализатора и графического редактора.
		Задачи:
		\begin{itemize}
			\item Разработать синтаксический анализатор файлов формата BMP (24-bit). Разобрать заголовок (File Header, DIB Header) и содержимое файла.
			\item Реализовать графические фильтры: инверсия, пороговое черно-белое преобразование, усиление резкости (матричная свертка), изменение яркости и контрастности, отражение по горизонтали и вертикали.
			\item Реализовать эффект «Глитч» (случайное смещение участков изображения).
			\item Реализовать функцию рирпроекции (хромакей): наложение изображения на фон с использованием маски по цвету (Green Screen / Blue Screen) с настраиваемым допуском.
			\item Организовать консольный пользовательский интерфейс (CLI) для взаимодействия с программой.
		\end{itemize}
	\end{enumerate}
	
	\newpage
	\section{Особенности реализации}
	Согласно заданию, решение разбито на два независимых проекта \texttt{stack}. Для работы с бинарными данными и обеспечения высокой производительности использовались библиотеки \texttt{bytestring}, \texttt{binary} и \texttt{vector}.
	
	\subsection{Часть 1: Реализация по книге Ф. Хагенлохера}
	
	В соответствии с главой 12 книги, для разбора текстовых и бинарных заголовков формата PNM была выбрана библиотека \texttt{attoparsec}, так как она обеспечивает высокую производительность и удобный интерфейс комбинаторов.
	
	В соответствии с главой 13, для представления данных использован подход GADT (Generalized Algebraic Data Types). Это позволяет определить единый тип \texttt{PNM}, конструкторы которого (\texttt{Greyscale} и \texttt{RGB}) явно кодируют формат данных, предотвращая ошибки типа ``Return-type polymorphism''.
	
	Для демонстрации возможностей многопоточности Haskell применена библиотека \texttt{parallel}. Реализована функция \texttt{processImageParallel}, которая разбивает вектор пикселей на фрагменты (чанки) и обрабатывает их параллельно на доступных ядрах процессора, используя стратегию \texttt{parList rdeepseq}.
	
	\subsection{Часть 2: Модификатор изображений BMP}
	
	\subsubsection{Типы данных}
	Для хранения изображения используется структура \texttt{Image}, содержащая ширину, высоту и плоский вектор байтов (\texttt{Word8}). Использование \texttt{Data.Vector.Storable} обеспечивает доступ к пикселям за время $O(1)$ и эффективное использование памяти, в отличие от списков.
	
	\begin{lstlisting}[caption={Определение типа Image}, label={lst:image_type}]
		module Types where
		
		import qualified Data.Vector.Storable as V
		import Data.Word (Word8)
		
		data Image = Image {
			imgWidth  :: Int,
			imgHeight :: Int,
			imgData   :: V.Vector Word8 
		} deriving (Show, Eq)
		
		data RGB = RGB { r :: Word8, g :: Word8, b :: Word8 } 
		deriving (Show, Eq)
	\end{lstlisting}
	
	\subsubsection{Синтаксический анализ BMP}
	Реализован собственный парсер формата BMP с использованием монады \texttt{Get} из библиотеки \texttt{Data.Binary}.
	
	Особенности реализации парсера:
	\begin{itemize}
		\item Чтение сигнатуры ``BM'' и заголовков (File Header, DIB Header).
		\item Проверка битности (поддержка 24-bit TrueColor).
		\item Учет выравнивания строк (Padding): каждая строка в BMP дополняется нулями до длины, кратной 4 байтам.
		\item Обработка порядка строк: BMP хранит данные снизу-вверх, поэтому при чтении строки разворачиваются.
	\end{itemize}
	
	\begin{lstlisting}[caption={Фрагмент парсера BMP}, label={lst:bmp_parser}]
		parseBMP :: Get Image
		parseBMP = do
		sig <- getWord16le -- "BM"
		skip 8 
		offset <- getWord32le
		headerSize <- getWord32le
		w <- getWord32le
		h <- getWord32le

		
		let rowSize = widthInt * 3
		let padding = (4 - (rowSize `mod` 4)) `mod` 4
		
		rawRows <- replicateM heightInt $ do
		row <- getByteString rowSize
		skip padding 
		return row
		
		let flatData = BS.concat (reverse rawRows) 
		return $ Image widthInt heightInt (V.fromList $ BS.unpack flatData)
	\end{lstlisting}    
	
	\subsubsection{Реализация фильтров}
	В модуле \texttt{Filters.hs} реализованы алгоритмы обработки изображений.
	
	\textbf{Пороговое черно-белое преобразование (Threshold):}
	Сравнение значения каждого канала с заданным порогом.
	
	\begin{lstlisting}[caption={Реализация порогового преобразования}, label={lst:bw}]
		thresholdBW :: Int -> Image -> Image
		thresholdBW level (Image w h v) = Image w h (V.map check v)
		where
		lvl = fromIntegral level
		check x = if x > lvl then 255 else 0
	\end{lstlisting}
	
	\textbf{Усиление резкости (Sharpen):}
	Реализовано через свертку (convolution) с матрицей ядра:
	$$
	\begin{pmatrix}
		0 & -1 & 0 \\
		-1 & 5 & -1 \\
		0 & -1 & 0 
	\end{pmatrix}
	$$
	Для каждого пикселя вычисляется взвешенная сумма его соседей.
	
	\textbf{Эффект «Глитч» (Glitch):}
	Алгоритм случайным образом выбирает два прямоугольных блока изображения и меняет их местами, создавая эффект цифрового искажения. Используется монада \texttt{IO} для генерации случайных координат.
	
	\subsubsection{Рирпроекция (Chroma Key)}
	Реализована функция наложения переднего плана (\texttt{fg}) на фон (\texttt{bg}).
	Алгоритм проверяет Евклидово расстояние между цветом текущего пикселя и ключевым цветом (хромакеем). Если расстояние меньше заданного допуска (\texttt{tolerance}), пиксель берется из фона, иначе — из переднего плана.
	
	\begin{lstlisting}[caption={Логика хромакея}, label={lst:chroma}]
		isKey :: (Word8, Word8, Word8) -> RGB -> Int -> Bool
		isKey (b, g, r) (RGB kr kg kb) tol =
		let dist = sqrt (db*db + dg*dg + dr*dr)
		in dist < fromIntegral tol
		
		chromaKey :: Image -> Image -> RGB -> Int -> Image
		chromaKey fg bg keyColor tol = ...

	\end{lstlisting}
	
	\subsubsection{Взаимодействие с пользователем}
	Реализован цикл обработки команд (\texttt{REPL}) в модуле \texttt{Main.hs}. Программа хранит текущее состояние (\texttt{AppState}), содержащее загруженное изображение. Поддерживаются команды загрузки, сохранения, применения эффектов и генерации тестовых изображений.
	
	\newpage
	\section {Результаты работы программы}
	\subsection{Интерфейс программы}
	Взаимодействие осуществляется через консоль. Пользователь может загружать файлы, применять цепочки эффектов и сохранять результат.
	
	Пример сессии работы:
	\begin{verbatim}
		> gen 400 300
		Generating test gradient 400x300...
		Generated.
		> sharp
		Applying Sharpen filter...
		Sharpening applied.
		> save test_sharp.bmp
		Saving to test_sharp.bmp...
		Saved.
	\end{verbatim}
	
		\begin{figure}[h!]
		\centering
		
\includegraphics[width=0.5\linewidth]{/Users/Seva/Desktop/1.png}
		\caption{Исходное изображение}
		\label{fig:chroma}
	\end{figure}
	
	Реализована обработка ошибок: программа корректно сообщает о неверном формате файла, отсутствии файла или неверных аргументах.
	
	\newpage
	\subsection{Демонстрация эффектов}
	Ниже представлены результаты работы фильтров на тестовом изображении (градиент и фотографии).
	
	\begin{figure}[h!]
		\centering
		
\includegraphics[width=0.45\linewidth]{/Users/Seva/Desktop/2.png}
		
\includegraphics[width=0.45\linewidth]{/Users/Seva/Desktop/3.png}
		\caption{Слева: инверсия. Справа: пороговое Ч/Б.}
		\label{fig:effects}
	\end{figure}
	
	\begin{figure}[h!]
		\centering
		
\includegraphics[width=0.45\linewidth]{/Users/Seva/Desktop/4.png}
		
\includegraphics[width=0.45\linewidth]{/Users/Seva/Desktop/5.png}
	 	\caption{Слева: эффект Глитч. Справа: вертикальный поворот.}
	 	\label{fig2:effects}
	 \end{figure}
	 
	 \begin{figure}[h!]
	 	\centering
	 	
\includegraphics[width=0.45\linewidth]{/Users/Seva/Desktop/6.png}
	 	\includegraphics[width=0.45\linewidth]{/Users/Seva/Desktop/7.png}
	 	\caption{Слева: горизонтальный поворот. Справа: резкость.}
	 	\label{fig3:effects}
	 \end{figure}
	 
	 	\newpage
	 \mbox{}
	 
	\begin{figure}[h!]
		\centering
		
\includegraphics[width=0.45\linewidth]{/Users/Seva/Desktop/8.png}
		
\includegraphics[width=0.45\linewidth]{/Users/Seva/Desktop/9.png}
		\caption{Слева: контраст. Справа: хромакея.}
		\label{fig4:effects}
	\end{figure}
		
			На примерах видно как работает последовательные операции и что их эффекты накладываются друг на друга. 
		
	\newpage
	\subsection{Рирпроекция (Хромакей)}
	Реализован механизм замены фона. Пользователь может выбрать пресет цвета (green/blue) или задать RGB вручную.
	
	\begin{figure}[h!]
		\centering
		
\includegraphics[width=0.5\linewidth]{/Users/Seva/Desktop/1.png}
		\caption{Фон хромакей}
		\label{fig5:chroma}
	\end{figure}
	
	Пример команды:
	\texttt{chroma bg.bmp green 100} — накладывает текущее изображение на \texttt{bg.bmp}, удаляя зеленый цвет с допуском 100.
	
	\begin{figure}[h!]
		\centering
		
\includegraphics[width=0.45\linewidth]{/Users/Seva/Desktop/11.png}
		
\includegraphics[width=0.45\linewidth]{/Users/Seva/Desktop/12.png}
		\caption{Слева: изображение до. Справа: изображение после.}
		\label{fig6:effects}
	\end{figure}


	\newpage
	\section*{Заключение}
	\addcontentsline{toc}{section}{Заключение}
	В рамках курсовой работы были разработаны два проекта на языке Haskell для обработки растровой графики.
	
	В первой части были освоены принципы комбинаторного парсинга (Attoparsec), использование GADT для строгой типизации и методы параллельного программирования для ускорения обработки изображений формата PNM.
	
	Во второй части был создан полноценный графический модификатор для формата BMP. Ключевые результаты:
	\begin{itemize}
		\item Написан собственный парсер бинарного формата BMP, корректно обрабатывающий заголовки и выравнивание.
		\item Использованы эффективные структуры данных (\texttt{Data.Vector.Storable}) для работы с пикселями, что обеспечило высокую скорость работы фильтров.
		\item Реализован широкий спектр алгоритмов: от простой инверсии до матричной свертки (резкость) и композитинга (хромакей).
		\item Создан удобный консольный интерфейс с поддержкой генерации тестовых данных.
	\end{itemize}
	
	Полученные навыки работы с бинарными данными, монадами состояния и векторами демонстрируют возможности Haskell в области системного программирования и обработки медиаданных.
	
	\newpage
	\section*{Список литературы}
	\addcontentsline{toc}{section}{Список литературы}
	
	\vspace{-1.5cm}
	\begin{thebibliography}{0}
		\bibitem{hagenlocher}
		Philipp Hagenlocher. Haskell from the Very Beginning. -- Leanpub, 2020.
		
		\bibitem{bmp_spec}
		BMP File Format. URL: \url{https://en.wikipedia.org/wiki/BMP_file_format} (дата обращения: 15.05.2025).
		
		\bibitem{hackage}
		Hackage -- binary: Binary serialisation for Haskell values using lazy ByteStrings. URL: \url{https://hackage.haskell.org/package/binary}, (дата обращения: 15.05.2025).
		
		\bibitem{vector}
		Hackage -- vector: Efficient Arrays. URL: \url{https://hackage.haskell.org/package/vector}, (дата обращения: 15.05.2025).
	\end{thebibliography}
\end{document}